\documentclass[12pt,a4paper,oneside]{article}
\usepackage[brazil]{babel}
\usepackage[utf8]{inputenc}
\usepackage{olymp}
\usepackage{color}
\usepackage[dvipsnames]{xcolor}
\usepackage{listings}

\setcounter{problem}{3}

\begin{document}

\begin{problem}{Datas válidas}{data}{2}

Faça um programa que receba uma data e informe se ela é uma data válida ou inválida.

Para isso faça obrigatoriamente uma função com o seguinte protótipo:

\texttt{\textcolor{Mulberry}{bool} validaData(\textcolor{Mulberry}{int} d, \textcolor{Mulberry}{int} m, \textcolor{Mulberry}{int} a)}

A função tem 3 parâmetros inteiros, representando o dia, o mês e o ano de uma data, e retorna um valor lógico, sendo verdadeiro caso a data seja uma data válida ou falso caso contrário.


%\Notes
%
%Observações, se tiver.

\InputFile

A entrada contém apenas uma linha, com três números inteiros, que indicam respectivamente dia, mês e ano. Restrições: $1 \leq$ dia $\leq 31, 1 \leq$ mês $\leq 12, 1592 \leq$ ano $\leq 3999$.

\OutputFile

Seu programa deve gerar apenas uma linha de saída, contendo ``\texttt{data valida}'' quando a data informada na entrada for válida ou ``\texttt{data invalida}'' quando não for. 

% \Notes
% Nesta questão, seu programa deve usar a função \texttt{main} dada acima exatamente como está, não a modifique! Sua tarefa é apenas criar a função \texttt{eh\_primo}.

\Example

\begin{example}
\exmp{
26 05 2023
}{
data valida
}%
\end{example}

\begin{example}
\exmp{
31 04 1983
}{
data invalida
}%
\end{example}

\begin{example}
\exmp{
31 02 2001
}{
data invalida
}%
\end{example}

\begin{example}
\exmp{
29 02 2023
}{
data invalida
}%
\end{example}

\begin{example}
\exmp{
29 02 2020
}{
data valida
}%
\end{example}



%Comentários sobre o exemplo, se necessário.

\end{problem}

\end{document}
